\section{Fundamentação teórica}
\label{literatura}

\subsection{Incentivos econômicos e alternativas de uso da terra}

Sob a ótica econômica, a escolha entre conservar e converter depende dos benefícios e custos das alternativas acessíveis ao produtor. A receita agrícola é apropriável, enquanto parte dos serviços mantidos pela vegetação beneficia outros agentes, de modo que a alternativa privada mais atraente pode divergir daquela que considera as perdas ambientais \cite{Angelsen2010}. O custo de oportunidade da conservação corresponde ao benefício líquido da melhor alternativa a que se renuncia, e não apenas ao valor da produção que poderia ocupar a área. Essa comparação reúne condições agronômicas, custos de abertura e produção, acesso a mercados e restrições institucionais. A participação do ambiente no processo econômico aparece tanto nas possibilidades de produzir quanto nos serviços perdidos com a conversão.

A conservação tampouco se resume a renunciar à produção presente. Na perspectiva de \citeonline{Daly1996}, a atividade econômica depende de uma base natural finita; em \citeonline{Dasgupta2001}, a consideração do capital natural vincula a riqueza às possibilidades de bem-estar futuro. Para o problema deste artigo, essa conexão significa que as condições ambientais participam das oportunidades produtivas e podem ser modificadas pelo próprio uso da terra. O fundamento é mais amplo que o teste empírico: examinar conversão e desempenho agrícola não permite valorar todos os serviços perdidos nem identificar o efeito da mudança ambiental sobre a produtividade futura.

Os preços participam dessa escolha sem determinar uma resposta territorial uniforme. Mantidos os custos e as demais condições, uma alta da cotação representa maior acréscimo de receita por hectare onde a produção potencial da cultura é maior. Essa vantagem pode favorecer o aproveitamento de áreas abertas ou tornar mais atraente a expansão sobre vegetação. \citeonline{LubowskiEtAl2008} relacionaram decisões de uso da terra a retornos líquidos e qualidade da terra nos Estados Unidos, enquanto \citeonline{Pfaff1999} destacou características locais e custos de transporte na Amazônia. O potencial agronômico ajuda a caracterizar a oportunidade de receita, mas sua realização depende de investimento, acesso e manejo. Por isso, um indicador de capacidade sob condições padronizadas não representa, isoladamente, rentabilidade ou tecnologia efetivamente empregada.

Em escala tropical, \citeonline{BermanEtAl2023} combinaram cotações internacionais com aptidão de quinze culturas em células geográficas entre 2001 e 2018. Encontraram associação positiva entre seu índice de preços e perda florestal, com diferenças conforme condições locais de integração comercial e desenvolvimento econômico. O estudo oferece um antecedente direto para examinar a exposição diferencial a preços comuns. A questão que permanece para o recorte de Mato Grosso é como essa associação se apresenta ao distinguir o potencial de soja e milho, o desempenho agrícola anterior e o destino imediato da conversão. A comparação requer medidas próprias, pois um índice composto de incentivos e a perda de cobertura arbórea não equivalem ao preço de uma cultura e à transição de vegetação nativa para agricultura.

Também importa distinguir o uso que cresce daquele que cede espaço. Na bacia do Tapajós, \citeonline{Braganca2018} encontrou que a valorização das lavouras em relação à carne favoreceu a substituição de pastagens por cultivos e reduziu o desmatamento local, sob condições próprias de realocação de fatores. Esse resultado qualifica a expectativa de que um incentivo agrícola mais elevado se traduza necessariamente em abertura: a resposta depende das alternativas já disponíveis. Para interpretar a associação entre preços e conversão, é preciso reconhecer tanto a possibilidade de expandir sobre vegetação quanto a de reorganizar a produção em terras anteriormente abertas.

\subsection{Transformação produtiva, biomas e instituições em Mato Grosso}

Mato Grosso torna essa distinção particularmente relevante pela coexistência de expansão agrícola, intensificação e conversão para pastagem. Na porção florestal do estado, \citeonline{MacedoEtAl2012} observaram que o aumento da produção de soja entre 2001 e 2005 decorreu da expansão sobre pastagens e florestas. Entre 2006 e 2010, a expansão continuou importante, mas 91\% dela ocorreu em terras previamente abertas, enquanto novas pastagens responderam pela maior parte do desmatamento remanescente. O resultado mostra que crescimento da produção e abertura direta para lavouras podem seguir trajetórias diferentes. Também delimita sua validade: a experiência da porção florestal naquele período não descreve automaticamente todos os biomas nem todo o intervalo examinado neste artigo.

A sucessão soja--milho acrescenta outra margem de crescimento, porque permite obter duas colheitas em uma área já cultivada. Ao discutir as transformações de Mato Grosso e Goiás, \citeonline{Spera2017} relacionou a intensificação às condições agronômicas, à disponibilidade de terras e aos incentivos para utilizar áreas abertas. A ampliação da produção, nesse caso, pode resultar de maior frequência de cultivo, sem igual acréscimo de área física. Para o presente estudo, essa experiência justifica examinar os preços das duas culturas e distinguir produção de conversão. Ela também indica que potencial produtivo elevado não determina sozinho a direção da pressão sobre a vegetação: pode favorecer tanto o uso mais intenso da terra quanto sua expansão.

As diferenças ambientais alteram o conteúdo dessa escolha. Na Amazônia mato-grossense, a perda de formações florestais e a relação entre pastagens e lavouras situam o problema discutido por \citeonline{MacedoEtAl2012}. No Cerrado, a diversidade de formações florestais, savânicas e campestres exige que a conservação alcance vegetação que não se resume à floresta de dossel fechado \cite{Spera2017}. A escolha da conversão bruta de vegetação nativa como desfecho acompanha essa diversidade. Permite observar a abertura para um uso determinado sem compensá-la pela regeneração em outra parcela, embora não identifique, por si, os serviços ambientais perdidos ou a legalidade da transição.

No Pantanal, a relação entre atividade econômica e conservação envolve condições distintas das áreas de cultivo mecanizado. \citeonline{SeidlEtAl2001} examinaram incentivos à implantação de pastagens cultivadas para a pecuária e identificaram associações com riqueza agropecuária, esforço produtivo e localização em relação à infraestrutura e aos mercados. O estudo também situou a magnitude e a duração das inundações anuais entre as restrições ao uso das terras. Assim, uma medida padronizada de potencial para soja ou milho não resume as possibilidades econômicas de uma paisagem inundável. A comparação com pastagem é pertinente, mas exige reconhecer a diferença entre o potencial agrícola calculado e as atividades efetivamente praticadas. Como o painel deste artigo contém poucas observações elegíveis no Pantanal, sua presença territorial não assegura uma estimativa representativa da dinâmica do bioma.

Essas condições ambientais se articulam com regras de abrangência diferente. A Lei nº 12.651/2012 estabelece obrigações de proteção da vegetação por imóvel e diferencia os percentuais de Reserva Legal conforme a localização e a formação vegetal, sem que a classificação predominante de um município determine todas as obrigações de suas propriedades \cite{Brasil2012}. Ao analisar Mato Grosso, \citeonline{SticklerEtAl2013} examinaram padrões de cumprimento, custos privados e consequências ambientais das mudanças no Código Florestal. A análise evidenciou a importância de distinguir o benefício coletivo da conservação dos custos percebidos pelos proprietários. A cobertura florestal municipal pode contextualizar esse problema, mas não informa diretamente passivos, obrigações históricas ou cumprimento individual.

A Moratória da Soja, por sua vez, vinculou o acesso a mercados a critérios de origem da produção no bioma Amazônia. \citeonline{HeilmayrEtAl2020} estimaram redução do desmatamento em áreas aptas à soja após o acordo, utilizando comparações de tempo, bioma e aptidão. Sua abrangência territorial e comercial difere da legislação aplicável aos imóveis, razão pela qual uma mesma cotação pode ser acompanhada de alternativas distintas dentro do estado. A existência de uma regra tampouco equivale à sua implementação, distinção central na interpretação das instituições \cite{North1991}.

Além do nível de desmatamento, interessa saber se os instrumentos alteram a resposta aos incentivos. \citeonline{HardingEtAl2021} examinaram a lista de municípios prioritários, a Moratória e áreas protegidas na Amazônia brasileira. Encontraram respostas diferentes entre instrumentos, incluindo maior sensibilidade a preços de outras culturas nos municípios expostos à Moratória, resultado compatível com substituição produtiva. Esse antecedente impede tratar uma mudança na associação com preços como equivalente ao efeito médio de uma política. No presente artigo, as frações territoriais e os marcos temporais permitem examinar heterogeneidade, sem reproduzir aqueles contrafactuais ou atribuir a cada bioma uma trajetória institucional uniforme.

\subsection{Perdas anteriores, recursos e expansão territorial}

A capacidade potencial não descreve todas as condições em que uma localidade responde aos preços. O desempenho das safras anteriores pode alterar recursos e alternativas de investimento, inclusive onde as condições agronômicas são favoráveis. É necessário, porém, distinguir uma produção habitualmente baixa de uma quebra em relação ao desempenho que a própria localidade vinha apresentando. Também se deve separar a perda física do resultado financeiro: receitas dependem de quantidade e preço, enquanto a disponibilidade de recursos depende ainda dos custos, das dívidas e dos mecanismos de compensação. Essas distinções orientam o uso de uma referência produtiva histórica, separada da medida de potencial.

A evidência sobre financiamento e risco oferece razões para considerar respostas distintas após uma perda. Na Amazônia brasileira, \citeonline{AssuncaoEtAl2020Credito} estimaram redução do desmatamento após o condicionamento ambiental do crédito, com efeito mais evidente em municípios voltados à pecuária e estimativa imprecisa nos voltados às lavouras. Em Gana, \citeonline{KarlanEtAl2014} encontraram aumento do investimento agrícola com seguro, distinguindo risco não segurado de falta de capital. Esses resultados sustentam a relevância dos recursos e da proteção contra riscos para as decisões produtivas. Não demonstram, entretanto, a sequência específica pela qual uma quebra de soja ou milho em Mato Grosso reduziria financiamento e conversão posterior.

Perdas produtivas também foram acompanhadas de expansão territorial. Na Zâmbia, \citeonline{OrdonezEtAl2025} examinaram uma infestação de lagarta-do-cartucho e encontraram perdas, aumento da produção de carvão e expansão de lavouras. A expansão foi imprecisa no painel de famílias e apareceu na análise espacial, distinção que impede reunir extração e abertura numa única resposta. Em texto para discussão, \citeonline{HsiaoEtAl2026} relacionaram calor extremo, perda florestal e expansão de lavouras. Na análise brasileira, baseada em censos agropecuários, encontraram redução da produtividade física acompanhada de aumento dos preços agrícolas locais. Os dois estudos mostram que choques adversos podem ser seguidos de maior pressão territorial, mas suas medidas de choque, escalas e mecanismos não se confundem com perdas municipais anteriores combinadas a cotações internacionais.

Para Mato Grosso, portanto, a restrição de recursos e a expansão após uma perda constituem possibilidades explicativas concorrentes. Uma safra desfavorável pode limitar o investimento necessário para aproveitar preços elevados; também pode ser seguida de expansão quando houver recursos e alternativas para reorganizar a produção. A associação da perda com a conversão, isoladamente, não distingue essas possibilidades nem responde se a relação com os preços se modifica. Essa é a função da interação examinada em H1b. A precedência da safra organiza a comparação, mas não transforma o indicador em choque exógeno nem revela as expectativas e os motivos individuais dos produtores.

\subsection{Questões em aberto e contribuição do artigo}

Os estudos discutidos estabelecem antecedentes para cada parte do problema, mas não tornam intercambiáveis seus objetos. A combinação entre preços e aptidão já foi explorada por \citeonline{BermanEtAl2023}; a coexistência de expansão produtiva e menor desmatamento foi documentada em Mato Grosso por \citeonline{MacedoEtAl2012}; e a relação entre choques produtivos e abertura foi examinada por \citeonline{OrdonezEtAl2025} e \citeonline{HsiaoEtAl2026}. A questão específica deste artigo é se o desempenho anterior modifica a associação entre preços e conversão bruta para agricultura, quando as diferenças de capacidade potencial são consideradas no mesmo painel municipal.

Um antecedente recente e particularmente próximo é \citeonline{RichardsArima2026}. Os autores construíram medidas de lucros por hectare e do conjunto da produção de soja e milho de segunda safra e as relacionaram a preços da terra e perda florestal na Amazônia de Mato Grosso. Encontraram associações positivas entre lucros setoriais e abertura, inclusive com respostas distribuídas no tempo. Essa evidência aproxima resultados econômicos anteriores e transformação territorial, mas não corresponde ao contraste preço×perda física aqui examinado. O presente artigo não mede lucros nem testa valorização patrimonial. Sua contribuição deve ser avaliada pela informação que o desempenho produtivo municipal acrescenta à heterogeneidade da associação com preços, sem reivindicar a identificação daqueles canais financeiros.

Essa comparação conduz a três escolhas complementares. A distinção entre potencial e desempenho realizado permite separar condições produtivas relativamente persistentes de perdas anteriores. A observação das transições para agricultura e pastagem distingue destinos que uma medida agregada de expansão ou perda florestal pode reunir. Por fim, a comparação dos modelos nas mesmas observações permite avaliar quanto uma mudança de resultado acompanha a seleção de localidades com histórico suficiente, antes de atribuí-la à inclusão das perdas. A contribuição reside no que esses contrastes permitem aprender em conjunto, inclusive quando a evidência é imprecisa ou contrária à direção inicialmente proposta.

A importância dessa investigação para a gestão ambiental decorre da possibilidade de distinguir situações que poderiam parecer semelhantes sob uma leitura apenas dos preços ou da produção estadual. Municípios com potencial elevado podem ampliar a produção em áreas abertas; perdas anteriores podem coexistir com possibilidades diferentes de financiamento e expansão; e as regras podem alcançar atividades e parcelas distintas. Examinar essas condições ajuda a delimitar onde uma associação é observável e quais informações faltam para interpretar seus mecanismos. O desenho não produz uma avaliação causal dos instrumentos de conservação, mas oferece uma comparação municipal de incentivos, condições produtivas e destinos da abertura, com limites territoriais e temporais explícitos.
